\chapter{Compétences apportées et développées par l'alternant}

\section{Capitalisation sur la double formation apportée par l'INSA CVL}



\begin{table}[H]
    \centering
    \begin{tabular}{lccp{6cm}}
        \toprule
        \textbf{Matière} & \textbf{Spécialité} & \textbf{Semestre} & \textbf{Compétences apportées par le module} \\
        \midrule
        Projet Mathématiques & MRI & S5 & \\
        Thermique & MRI & S5 & \\
        Composants Électroniques & MRI & S5 & \\
        Optimisation Linéaire & MRI & S6 & \\
        Analyse Numérique & MRI & S6 & \\
        Automatique & MRI & S6 & \\
        Modélisation et commande dans l'espace d'état & MRI & S7 & \\
        Projet Informatique et Sécurité & STI & S7 & \\
        Introduction au Machine Learning & STI & S8 & \\
        Deep Learning et applications & STI & S8 & \\
        Traitement de données NoSQL & STI & S8 & \\
        Data Mining & STI & S9 & \\
        IA Avancée & STI & S9 & \\
        \bottomrule
    \end{tabular}
    \caption{Apports de la formation INSA pour l'alternance}
    \label{tab:modules}
\end{table}

\ldots

\begin{itemize}
    \item Par le stage de 4A MRI, portant sur le développement d'un outil de liquéfaction des barrages en remblai :
    \begin{itemize}
        \item
    \end{itemize}
    \item Par le stage de 4A STI, portant sur l'étude de l'architecture logicielle d'un code de calcul en acoustique environnementale pour les installations industrielles du parc EDF :
    \begin{itemize}
        \item
    \end{itemize}
\end{itemize}

\section{Marges de progression identifiées et continuité de la formation}



\section{Premier retour d'expérience et suite immédiate de l'alternance}